\section{Angular fields with square-zero differential}
\label{c:nil-section}

We now use the limiting field from the end of
Section~\ref{c:sing-edge-one}, which is constant along each ray of its
triangular label domain.  Its differential generates a shear, so
the square-zero case of the finite approximation theorem applies
away from the apex.  We prove closed coverage by two such sweeps, compute
their cost, and then cover the omitted apex segments by a finite net.

Return to the edge-two tetrahedron $K_2$ of Section~\ref{c:sing-fields},
using the normalized height $z$ in place of $t$ and setting $h=\sqrt2$:
\begin{equation}\label{c:nil-tetrahedron}
 K_2=
 \{(x,y,hz):0\le z\le1,\ |x|\le1-z,\ |y|\le z\},
 \qquad h=\sqrt2.
\end{equation}
Its minimum projection area is $h$, as established by similarity in
Section~\ref{c:sing-fields}.
Thus $(x,y,z)$ are normalized coordinates: their Euclidean metric is
$dx^2+dy^2+2\,dz^2$.  The physical similarity
\begin{equation}\label{c:nil-similarity}
 (x,y,hz)\longmapsto
 \left(\frac{x}{2},\frac{y}{2},\frac{1/2-z}{\sqrt2}\right)
\end{equation}
has ratio $1/2$ and takes $K_2$ to the centered edge-one
tetrahedron $K_1$: writing $t=1/2-z$, its image satisfies
$|X|\le(1/2+t)/2$, $|Y|\le(1/2-t)/2$, and has physical height
$t/\sqrt2$.  Both perpendicular base areas and projection areas are
divided by four.  We retain the present scale throughout this section.

\subsection{Choosing the angular field}

Write a nonzero label in the upper half-plane as $b=(pt,t)$,
where $t>0$.  For a field of the form $G(b)=(v(p),w(p))$,
differentiation gives
\[
 DG(b)=\frac1t
 \begin{pmatrix}v'(p)\\w'(p)\end{pmatrix}
 \begin{pmatrix}1&-p\end{pmatrix}.
\]
Thus $v'(p)=p w'(p)$ makes $(DG)^2=0$: the row in this product
annihilates its column.  We also want the velocity tangent to the
boundary rays $p=\pm1$, which requires $v(p)=p w(p)$ there.
Fix $c=1/2$ and choose $w(p)=cp$. Integrating the differential relation
with these endpoint conditions determines $v$. We also define $a$ by
\begin{equation}\label{c:nil-angular-functions}
 v(p)=\frac c2(1+p^2),\qquad w(p)=cp,
 \qquad a(p)=c^2p^2v(p).
\end{equation}
The last function specifies the quadratic overlap we will allow.
For $0<k<1$, put
\[
 H(p)=c+k^2a(p)+k^3,\qquad -1\le p\le1.
\]
The extra cubic term provides a strict margin in the coverage proof.
In normalized coordinates define two sets swept out by segments:
\begin{equation}\label{c:nil-sweeps}
 \begin{aligned}
 \mathcal E_+&=\{(s,pt+ksv(p),t+ksw(p)):
     |s|\le1,\ |p|\le1,\ 0\le t\le H(p)\},\\
 \mathcal E_-&=\{(pt-ksv(p),s,1-t+ksw(p)):
     |s|\le1,\ |p|\le1,\ 0\le t\le H(p)\}.
 \end{aligned}
\end{equation}
At $k=0$ the two sweeps cover the lower and upper halves of the
tetrahedron by axis-parallel segments.  For positive $k$ we must check that
their varying tilts do not leave a gap between them.

\begin{lemma}\label{c:nil-sweep-cover}
For every $0<k<1$, every normalized point of
$K_2$ belongs to $\mathcal E_+\cup\mathcal E_-$.
\end{lemma}

\begin{proof}
We first record an elementary inversion including the endpoints.
If $z\ge0$, $d\in\R$ and $|y|\le z$, the equation
\[
 y=rz+d(1-r^2)
\]
has a solution $r\in[-1,1]$ with $z-2dr\ge0$.
For $d>0$, the nondecreasing part of this concave quadratic starts
at $r=-1$ with value $-z$ and reaches a value at least $z$.
For $d<0$, the nondecreasing part of the convex quadratic ends
at $r=1$ with value $z$ and starts at a value at most $-z$.
For $d=0$ the equation is linear; if also $z=0$, then $y=0$
and any $r$ is allowed.

For a normalized point $(x,y,z)$ of $K_2$, apply
this observation twice to obtain $r,q\in[-1,1]$ such that
\begin{equation}\label{c:nil-inversion}
 \begin{aligned}
 y&=rz+\frac{kcx}{2}(1-r^2),
 &t&=z-kcxr\ge0,\\
 x&=q(1-z)-\frac{kcy}{2}(1-q^2),
 &t'&=1-z+kcyq\ge0.
 \end{aligned}
\end{equation}
Equivalently,
\begin{equation}\label{c:nil-inverse-sweeps}
 y=rt+kxv(r),\qquad x=qt'-kyv(q).
\end{equation}
If $t\le H(r)$, the first sweep contains the point with axial
parameter $s=x$.  If $t'\le H(q)$, the second contains it with
$s=y$.  In both cases $|s|\le1$.

Suppose both inequalities fail.  Then $t,t'>c$, while
\[
 t+t'=1+kc(yq-xr)\le1+2kc.
\]
Since $2c=1$, both $t-c$ and $t'-c$ lie between zero and $2kc$.
The bounds $|x|,|y|\le1$ and $0<v\le c$, together with
\eqref{c:nil-inverse-sweeps}, now imply
\begin{equation}\label{c:nil-center-error}
 |x-cq|\le3kc,\qquad |y-cr|\le3kc.
\end{equation}
Substituting \eqref{c:nil-inverse-sweeps} into the expression for
$t+t'-1$ gives the exact identity
\begin{equation}\label{c:nil-overlap-identity}
 (1-kcrq)(t-c)+(1+kcrq)(t'-c)
 =k^2c\bigl(xqv(r)+yrv(q)\bigr).
\end{equation}
Both coefficients on the left are positive.  The assumed cutoff
failures make that side strictly greater than
\[
 k^2(a(r)+a(q))+
 \bigl(2+crq(a(q)-a(r))\bigr)k^3
 \ge k^2(a(r)+a(q))+(2-c^4)k^3,
\]
since $0\le a\le c^3$.  On the other hand,
\eqref{c:nil-center-error} bounds the right side by
\[
 k^2c^2\bigl(q^2v(r)+r^2v(q)\bigr)+6c^3k^3
 \le k^2(a(r)+a(q))+6c^3k^3.
\]
The last inequality follows from
$(r^2-q^2)(v(r)-v(q))\ge0$.
These bounds contradict $6c^3<2-c^4$ at $c=1/2$.
All inversions used closed intervals, so the same argument includes
the faces, edges and vertices.
\end{proof}

\subsection{The cost predicted by the sweeps}

We have proved complete coverage by two continuous families.  We next
compute their limiting cost, so that a definite positive tolerance is
available when making the cover finite.

For a planar tile $T$ and a constant velocity $g=(g_1,g_2)$, write
\[
 \mathcal C(T,g)=\{(s,\xi):s\in\R,\ \xi\in T+sg\},
 \qquad J_h(g)=(1+g_1^2+h^2g_2^2)^{-1/2}.
\]
The two physical maps
\begin{equation}\label{c:nil-physical-maps}
 F_+(s,\xi)=(s,\xi_1,h\xi_2),\qquad
 F_-(s,\xi)=(\xi_1,s,h(1-\xi_2))
\end{equation}
send this auxiliary cylinder to cylinders with perpendicular base
area $h|T|J_h(g)$.  Indeed, the transverse tile has physical area
$h|T|$, and projection to the plane perpendicular to
$(1,g_1,hg_2)$ has area factor $J_h(g)$.  The second map differs
from the first by a Euclidean isometry.  A square $T$, of any
orientation, therefore gives a compact convex parallelogram base.

Under these maps the two sweeps have the common auxiliary form
\begin{equation}\label{c:nil-signed-sweep}
 (s,(pt,t)+s\sigma k(v(p),w(p))),\qquad
 |s|,|p|\le1,\quad 0\le t\le H(p),
 \qquad \sigma\in\{1,-1\}.
\end{equation}
The Jacobian of $(p,t)\mapsto(pt,t)$ is $t$.  The integral
predicting the weighted tile cost of either sweep is consequently
\begin{equation}\label{c:nil-integral}
 I(k)=\int_{-1}^1\frac{H(p)^2}{2}
 \bigl(1+k^2(v(p)^2+h^2w(p)^2)\bigr)^{-1/2}\,dp.
\end{equation}

\begin{lemma}\label{c:nil-saving}
As $k\downarrow0$,
\[
 I(k)=\frac14-\frac{k^2}{480}+O(k^3).
\]
In particular $I(k)<1/4$ for every sufficiently small positive $k$.
\end{lemma}

\begin{proof}
All functions of $p$ in \eqref{c:nil-integral} are uniformly bounded
on $[-1,1]$.  Expanding the two factors, uniformly on that interval,
gives
\begin{equation}\label{c:nil-expansion}
 I(k)=c^2+k^2\left(
 c\int_{-1}^1a(p)\,dp-
 \frac{c^2}{4}\int_{-1}^1(v(p)^2+h^2w(p)^2)\,dp
 \right)+O(k^3).
\end{equation}
Direct polynomial integration yields
\[
 \int_{-1}^1a(p)\,dp=\frac{8c^3}{15},\qquad
 \int_{-1}^1(v(p)^2+2w(p)^2)\,dp=\frac{34c^2}{15}.
\]
The coefficient of $k^2$ in \eqref{c:nil-expansion} is therefore
$-c^4/30=-1/480$, as asserted.
\end{proof}

\subsection{A finite cover, including the apex}

We now realize the saving by finitely many cylinders.  On the open
half-plane $b_2>0$ define
\[
 G_\sigma(b)=\sigma k
 \bigl(v(b_1/b_2),w(b_1/b_2)\bigr),
\]
using the same polynomial formulas for $v,w$ at every real argument.
This is a smooth field, and the earlier differential computation gives
\begin{equation}\label{c:nil-rank-one}
 DG_\sigma(pt,t)=\frac{\sigma kc}{t}
 \begin{pmatrix}p\\1\end{pmatrix}
 \begin{pmatrix}1&-p\end{pmatrix},
 \qquad (DG_\sigma)^2=0.
\end{equation}
At $t=0$ every label $(pt,t)$ is the origin, but its assigned
velocity still depends on $p$.  Thus the field has no continuous
extension to the apex.  That part of the family must be covered
separately.

\begin{lemma}[A finite apex net]\label{c:nil-apex-net}
Fix $0<k<1$ and $0<\rho\le1$.  For either sign $\sigma$, the
part of \eqref{c:nil-signed-sweep} with $0\le t\le\rho$ is
covered by at most $\lceil2/\rho\rceil+1$ auxiliary cylinders.
Their tiles are squares of side $2(1+kc)\rho$, and their total
weighted tile area is at most
\begin{equation}\label{c:nil-apex-cost}
 16(1+kc)^2\rho.
\end{equation}
\end{lemma}

\begin{proof}
Put $m=\lceil2/\rho\rceil$ and $p_j=-1+2j/m$ for
$0\le j\le m$.  Every $p\in[-1,1]$ is within $\rho$ of one
$p_j$.  Use the common tile
$T_j=[-(1+kc)\rho,(1+kc)\rho]^2$ and velocity
$g_j=\sigma k(v(p_j),w(p_j))$.
Since $|v'|\le c$ and $|w'|=c$ on this interval, each coordinate
of
\[
 (pt,t)+s\sigma k\bigl((v(p),w(p))-(v(p_j),w(p_j))\bigr)
\]
has absolute value at most $(1+kc)\rho$ for
$0\le t\le\rho$ and $|s|\le1$.  The cylinders therefore
cover these segments, also at $t=0$.  Finally $J_h\le1$ and
$m+1\le2/\rho+2\le4/\rho$, so their weighted cost is at most
$(4/\rho)\,4(1+kc)^2\rho^2$.
\end{proof}

The remaining finite approximation uses the same square-tile theorem,
\BNilpotent{}, with $L=1$. A square-zero planar matrix has trace and
determinant zero, so it satisfies the theorem's condition
$\operatorname{tr}DG=0$ and $-1<\det DG\le0$.
For a fixed compact set $D$, a $C^1$ field $G$ on a neighborhood of $D$,
fixed $h>0$, and positive tolerance, the theorem supplies finitely
many compact nondegenerate square tiles, with possibly different
orientations, and velocities covering
$\{(s,b+sG(b)):b\in D,\ |s|\le1\}$.
Their weighted cost is at most $\int_D J_h(G)$ plus that tolerance.
The weights use the original velocity coordinates. Neither a
small-shear condition nor a boundary regularity assumption is required.

\begin{proposition}\label{c:nil-finite-counterexample}
The tetrahedron $K_2$ admits a finite cylinder cover
with compact convex parallelogram perpendicular bases whose total
area is less than
$\frac12\min_{|u|=1}|\pi_{u^\perp}K_2|$.
\end{proposition}

\begin{proof}
Choose and fix $0<k<1$ with $I(k)<c^2$, and then choose
$0<\epsilon<c^2-I(k)$.  Next fix
\[
 0<\rho<\min\left\{c,\frac{\epsilon}{32(1+kc)^2}\right\}.
\]
The apex net of Lemma~\ref{c:nil-apex-net} costs less than
$\epsilon/2$ in weighted tile area, for each sign.
Only now form the truncated label set
\[
 D_\rho=\{(pt,t):-1\le p\le1,\ \rho\le t\le H(p)\}.
\]
It is compact, lies in the open half-plane $b_2>0$, and the field
$G_\sigma$ is $C^1$ on that neighborhood.  Equation
\eqref{c:nil-rank-one} verifies its square-zero differential there,
hence trace and determinant zero.
We may therefore apply \BNilpotent{}
with this fixed $D_\rho$, physical factor $h=\sqrt2$, unit axial
interval, and tolerance $\epsilon/2$.  Although the derivative
may become large when $\rho$ is small, no hypothesis bounds it by
a prescribed constant; the finite approximation is chosen after
$\rho$ has been fixed.

Since the weight is independent of $t$ and of the sign,
\[
 \int_{D_\rho}J_h(G_\sigma(b))\,db
 =\int_{-1}^1\frac{H(p)^2-\rho^2}{2}
   \bigl(1+k^2(v(p)^2+h^2w(p)^2)\bigr)^{-1/2}\,dp
 \le I(k).
\]
Combining the approximation with the apex net gives a finite
physical cylinder cover of each entire sweep at cost at most
$h(I(k)+\epsilon)$.  All its perpendicular bases are compact
convex parallelograms by \eqref{c:nil-physical-maps}.
Lemma~\ref{c:nil-sweep-cover} now proves that the two lists
together cover the whole tetrahedron, and their total area is at most
\[
 2h(I(k)+\epsilon)<2hc^2=\frac h2
 =\frac12\min_{|u|=1}|\pi_{u^\perp}K_2|.
\]
The explicit apex count controls the omitted endpoint.  Finiteness
of the other lists is supplied at the fixed tolerance; no uniform
bound on their number is needed.
\end{proof}
